% ---------------------------------------------------------------------
%  TEORETICKÁ ČASŤ – max. 3 strany.
%  Podľa podmienok predmetu je v tejto časti ZAKÁZANÉ používať text generovaný
%  umelou inteligenciou (porušenie = 0 bodov). Text preto píšu autori sami.
%  Nižšie je iba navrhovaná osnova a zoznam overených zdrojov ku každej časti.
% ---------------------------------------------------------------------
\chapter{Analýza problematiky}

\todo{Celú kapitolu napíšte sami (max.\ 3 strany). Pri každej podkapitole sú uvedené
zdroje z \texttt{references.bib}, ktoré sú overené a nie staršie ako 5 rokov.}

\section{Biometrický systém rozpoznávania tváre}
\todo{Princíp: detekcia tváre, zarovnanie, extrakcia príznakov, porovnanie, rozhodnutie;
verifikácia (1:1) vs.\ identifikácia (1:N); hlbokové príznaky s uhlovou rezervou.
Zdroje: \cite{wu2023yunet, deng2022arcface, kim2022adaface, zhong2021sface}.}

\section{Kvalita biometrickej vzorky}
\todo{Prečo kvalita ovplyvňuje FRR; čo meria norma ISO/IEC 29794-5.
Zdroje: \cite{iso29794-5, schlett2022fiqa}.}

\section{Prezentačné útoky a ich detekcia}
\todo{Definícia prezentačného útoku a PAI (tlač, replay, maska); pasívne vs.\ aktívne
metódy detekcie živosti; problém generalizácie na nové typy útokov;
injekčné útoky ako hranica PAD.
Zdroje: \cite{yu2023fassurvey, hernandezortega2023pad, sun2023rethinking,
srivatsan2023flip, wang2023wildfas, enisa2022remote}.}

\section{Metriky a normy hodnotenia}
\todo{FAR/FRR/EER pre rozpoznávanie; APCER, BPCER, ACER podľa ISO/IEC 30107-3;
stručne prehľad noriem. Zdroje: \cite{iso30107-3, busch2023standards}.}

\section{Ochrana biometrických údajov}
\todo{Biometrické údaje ako osobitná kategória osobných údajov, ochrana šablón,
regulácia biometrických systémov v EÚ. Zdroje: \cite{abdullahi2024template, euaiact2024}.}
